\documentclass[english,twoside, openright, hidelinks, 11 pt, class=report,crop=false]{standalone}
\input{../../preamb}
\input{../bm}
\begin{document}

\section{Følger}\index{følge}
Følger er en oppramsing av tall i en bestemt rekkefølge, gjerne skilt med komma. I følgen 
\begin{equation}
2, 4, 8, 16  \label{folg}
\end{equation}
sier vi at vi har fire \outl{element}\index{element}. Element nr.\,1 har verdi 2, element nr.\,2 har verdi 4 og så videre. Hvert element i en følge beskrives gjerne ved hjelp av en indeksert bokstav. Velger vi oss bokstaven $ a $ for følgen over, kan vi skrive
\[ a_1=2\quad,\quad a_2=4\quad,\quad a_3=8\quad,\quad a_4=16 \] 

Når vi lar $ a_i $ betegne elementene i en følge, bruker vi $ {i\in \mathbb{N}}$. I likhet med mengder, kan vi bruke \sym{$ \lbrace\rbrace $} for å indikere en følge, og \sym{$ \in $} for å vise at et element er inneholdt i en følge. For eksempel er $ 8\in\lbrace2, 4, 8, 16\rbrace $ og $\lbrace a_i\rbrace=\lbrace2, 4, 8, 16\rbrace$. \vsk

Ofte kan tallene i en følge settes i sammenheng med hverandre. Multipliserer vi for eksempel et element i følgen fra \eqref{folg} med $ 2 $, så har vi funnet det neste elementet. Den \outl{rekursive formelen}\index{rekursiv formel} er da
\[ a_i = 2\cdot a_{i-1} \]
I den rekursive formelen bruker vi altså den forrige\footnote{Det kan også være nødvendig å bruke flere foregående verdier.} verdien for å finne den neste. \vsk

Av uttrykket for $a_i$ kan vi legge til uendelig mange nye element i \eqref{folg}, og får da følgen
\begin{equation}
 2, 4, 8, 16, 32, 64, ...  \label{folg1}
\end{equation}
Symbolet \sym{$ ... $} viser at nye element fortsetter i det uendelige.\vsk

Hva om vi for denne følgen ønsker å finne element nr. 20, altså $ a_{20} $? Det vil da lønne seg å finne en \outl{eksplisitt formel}\index{eksplisitt formel}. For å gjøre dette skriver vi opp noen element, og ser om vi finner et mønster: 
\alg{
& a_1 = 2 = 2^1 \\
& a_2 = 4 = 2^2 \\
& a_3 = 8 = 2^3
}
Av ligningene over innser vi at vi for element nr.\,$ i $ kan skrive
\[ a_i=2^i \] 
Og slik kan vi fort finne element nr.\,20:
\[ a_{20}=2^{20}= 1048576 \]
En eksplisitt formel gir oss et uttrykk der verdien til et element regnes ut. Når man har et slikt uttrykk er det også vanlig å skrive dette som siste element i rekken, \eqref{folg1} blir da seende slik ut:
\[  2, 4, 8, 16, 32, 64, ..., 2^i \]
\spr{
En følge med et konkret antall element kalles en \outl{endelig følge}\index{følge!endelig}. En følge med uendelig mange element kalles en \outl{uendelig følge}\index{følge!uendelig}. Følgen fra \eqref{folg} er altså en endelig følge, mens følgen fra \eqref{folg1} er en uendelig følge.\vsk

Om ''element nr. $i$'' i en følge sier vi gjerne det ''$i$-te elementet'' i følgen.
}
\info{Forskjell på følge og mengde}{
I en følge har rekkefølgen på tallene noe å si. Hvis $\lbrace{2, 4, 6\rbrace}$ og ${\lbrace2, 6, 4\rbrace}$ beskriver to følger, er ikke følgene like. Hvis de to uttrykkene derimot beskriver mengder, er de like.\vsk

En mengde inneholder ikke flere element med samme verdi, mens en følge kan gjøre det.
}

\subsection{Aritmetiske følger} \index{følge!aritmetisk}

Følgen 
\[ 2, 5, 8, 11, 14, 17 \]
kalles en \textit{aritmetisk følge}. I en aritmetisk følge har to naboelement konstant differanse, i dette tilfellet 3. Skriver vi opp de tre første elementene kan vi finne mønsteret til en eksplisitt formel\index{eksplisitt formel}:
\alg{
& a_1 = 2 = 2+3\cdot0\\
& a_2 = 5 = 2+3\cdot1\\
& a_3 = 8 = 2 +3\cdot2
}
Av ligningene over observerer vi at
\[ a_i=2+3(i-1) \]
\reg[Aritmetisk følge]{
	Et element $ a_i $  i en aritmetisk følge er gitt ved den rekursive formelen
	\begin{equation}\label{arrekeq}
		a_i=a_{i-1}+d
	\end{equation}
	og den eksplisitte formelen
	\begin{equation}
		a_i=a_1+d(i-1) \label{ekspl}
	\end{equation}
	hvor $ d $ er den konstante differansen $ a_i-a_{i-1} $.}
\newpage
\eks[]{Finn den rekursive og den eksplisitte formelen til følgen
	\[ 7, 13, 19, 25, ... \]	\vs
	\sv
	Følgen har konstant differanse $ {d=6} $ og første element $ {a_1=7} $. Den rekursive formelen blir da
	\[ a_i = a_{i-1}+6 \]
	Mens den eksplisitte formelen blir
	\[ a_i=7+6(i-1) \]
}
\subsection{Geometriske følger} \index{følge!geometrisk}
Følgen 
\[ 2, 6, 18, 54, 162 \]
kalles en \outl{geometrisk følge}. I en geometrisk følge har forholdet mellom to naboelement den samme kvotienten, i dette tilfellet 3. Også her kan vi gjenkjenne et fast mønster:
\algv{
	& a_1 = 2 = 2\cdot3^0\\
	& a_2 = 6 = 2\cdot3^1\\
	& a_3 = 18 = 2\cdot3^2
}
Den eksplisitte formelen blir derfor
\[ a_i = 2\cdot3^{i-1} \]
\reg[Geometrisk følge]{
	Et element $ a_i $  i en geometrisk følge med kvotient $ k $ er gitt ved den rekursive formelen
	\begin{equation}\label{geofolgrek}
		a_i=a_{i-1}\cdot k
	\end{equation}
	og den eksplisitte formelen
	\begin{equation}\label{geofolgeks}
		a_i=a_1\cdot k^{i-1}
	\end{equation}
}
\eks[]{
Finn den rekursive og den eksplisitte formelen til følgen
\[ 5, 10, 20, 40, ... \]

\sv
Følgen har konstant kvotient $ {k=2} $ og første element $ {a_1=5} $. Den rekursive formelen blir da
\[ a_i= a_{i-1}\cdot 2 \]
Mens den eksplisitte formelen blir
\[ a_i=5\cdot2^{i-1} \]
}


\section{Rekker}\index{rekke}
Ved å addere tall, danner vi en \outl{rekke}. Et eksempel er
\[ 2+6+18+54+162 \]
Vi indekserer leddene på samme måte som vi gjør for element i en følge; i rekken over har ledd nr. 3 verdien 18.\vsk

For en rekke er det naturlig at vi ikke bare ønsker å vite verdien til hvert enkelt ledd, men også hva summen av alle leddene er.
\spr{
En rekke med et konkret antall ledd kaller vi en \outl{endelig rekke}. En rekke med uendelig mange ledd kaller vi en \outl{uendelig rekke}. I mange matematiske tekster blir en rekke definert som addisjon mellom uendelig mange ledd. Det vi kaller en \textit{endelig rekke} beskrives da som en \outl{delsum} av en rekke. \vsk

Hvis summen av leddene til en uendelig rekke går mot en konkret verdi, sier vi at rekken er \outl{konvergent}\index{rekk!konverget} og at rekken \outl{konvergerer} mot denne verdien. I omvendt tilfelle sier vi at rekken er \outl{divergent}\index{rekke!divergent}.
}
\info{Merk}{
En rekke som
\[ 1+\left(-\frac{1}{3}\right)+\frac{1}{5}+\left(-\frac{1}{7}\right)+\frac{1}{9}+... \]
Skrives ofte bare som
\[ 1-\frac{1}{3}+\frac{1}{5}-\frac{1}{7}+\frac{1}{9}+... \]
}

\subsection{Aritmetiske rekker}
\reg[\sumarrek \label{sumarrek}]{
	
	\index{rekke!aritmetisk}
	Hvis leddene i en rekke kan beskrives som en aritmetisk følge, kalles rekken en \outl{aritmetisk rekke}.\vsk
	
	Summen $ S_n $ av de $ n $ første leddene i en aritmetisk rekke er gitt som
	\begin{equation}\label{arrek}
		S_n = n\frac{a_1+a_n}{2}
	\end{equation}
	hvor $ a_1 $ er det første elementet i rekken.
} \regv

\eks{ \label{arrekeks}
	Gitt den uendelige rekken
	\[ 3+7+11+... \] 
	Finn summen av de ti første leddene \os
		
	\sv
	Det $ i $-te leddet $a_i$ i rekken er gitt ved formelen
	\[ a_i=3+4(i-1) \]
	Dette er derfor en aritmetisk rekke, og summen av de $ n $ første  leddene er da gitt av \eqref{arrek}. Ledd nr.\,10 blir da
	\[ a_{10} = 3+4(10-1) = 39 \]
	Altså har vi at
	\[ S_{10} = 10\cdot\frac{3+39}{2} = 210 \]
} \vsk
\fork{\ref{sumarrek} \sumarrek}{
	Ved å bruke den eksplisitte formelen fra \eqref{ekspl}, kan vi skrive summen av en aritmetisk rekke med $ n $ element som
	\begin{equation}
		S_n = a_1 + (a_1+d) + (a_1+ 2d)+...+ (a_1+d(n-1)) \label{a1}
	\end{equation}
	Men et ledd i rekken kan også uttrykkes slik:
	\[ a_i= a_n-(n-i)d \]
	Og da kan vi skrive summen som (her står ledd nr. $n$ først, deretter ledd nr. $n-1$ osv.)
	\begin{equation}
		S_n = a_n + (a_n-d)+(a_n-2d)+...+(a_n-d(n-1)) \label{a2}
	\end{equation}
	Adderer vi \eqref{a1} og \eqref{a2}, får vi at
	\alg{
		2S_n &= na_1 + na_n \br
		S_n&=n\frac{a_1+a_n}{2} 
	}
}
\newpage
\subsection{Geometriske rekker} 
\reg[\sumgerek \label{geom} \index{rekke!geometrisk}]{
Hvis leddene i en rekke kan beskrives som en geometrisk følge, kalles rekken en \outl{geometrisk rekke}.	\vsk

Summen $ S_n $ av de $ n $ første leddene i en geometrisk rekke med kvotient $ k $ og første ledd $ a_1 $ er gitt som
	\begin{equation}
		S_n=a_1\frac{1-k^n}{1-k}\quad, \quad k\neq 1 \label{sumg}
	\end{equation}
	Hvis $ {k=1}, $ er
	\begin{equation}
		 S_n = na_1
	\end{equation}
}
\eks{ \label{gerekeks}
Gitt den uendelige rekken
\[ 3+6+12+24+... \]	
Finn summen av de 15 første leddene.

\sv
Dette er en geometrisk rekke med $ {a_1=3} $ og $ {k=2} $. Summen av de 15 første leddene blir da
\alg{
	S_{15}&= 3\cdot\frac{1-2^{15}}{1-2} \br
	&= 	3\cdot\frac{1-32768}{-1} \br
	&= 98301
	}
}
\newpage
\fork{\ref{geom} \sumgerek}{
Summen $ S_n $ av en geometrisk rekke med $ n $ element er
\begin{equation}
	S_n = a_1 + a_1k + a_1k^2+...+a_1 k^{n-2}+a_1 k^{n-1} \label{geo1}
\end{equation}
Ganger vi denne summen med $ k $, får vi at
\begin{equation}
	kS_n = a_1k + a_1k^2 + a_1k^3+...+a_1 k^{n-1}+a_1 k^{n} \label{geo2}
\end{equation}
\eqref{geo2} subtrahert \eqref{geo1} gir
\alg{
	S_n-kS_n &= a_1-a_1k^n \\
	S_n(1-k) &= a_1(1-k^n) \\
	S_n &= a_1\frac{1-k^n}{1-k}
}
}
\subsection{Uendelige geometrisk rekker}
For en geometrisk rekke med uendelig mange ledd merker vi oss dette: \vsk

Hvis $ {|k|<1} $, er
\[
\lim\limits_{n\to\infty} S_n =\lim\limits_{n\to\infty} a_1\frac{1-k^n}{1-k} 
= a_1 \frac{1}{1-k}\]
Summen av uendelig mange ledd går altså mot en konkret verdi. I dette tilfellet er altså rekken konvergent. Hvis derimot $ |k|\geq 1$, går summen mot $ \pm \infty$. Da er rekken divergent.\regv
\reg[Summen av en uendelig geometrisk rekke]{
	For en uendelig geometrisk rekke med kvotient $ {k<|1|} $ og første element $ a_1 $ er summen $ S_\infty $ av rekken gitt som
	\begin{equation}
		S_\infty = \frac{a_1}{1-k} \label{infsum}
	\end{equation}
}
\newpage
\eks{
Finn summen av rekken
\[ 7+7\cdot10^{-1}+7\cdot10^{-2}+... \]
\sv
Dette er en geometrisk rekke med $a_1=7$ og $k=10^{-1}=\frac{1}{10}$. Dermed har vi at
\[ S_\infty=\frac{7}{1-\frac{1}{10}}=\frac{70}{9} \]
}
\newpage
\subsection{Summetegnet}\index{summetegnet}
Vi skal nå se på et symbol som forenkler skrivemåten av rekker betraktelig. Symbolet blir spesielt viktig i \refkap{Integrasjon}.\vsk

Tidligere har vi skrevet rekkene mer eller mindre bent frem. For eksempel har vi sett på rekken
\[ 2+6+18+54+162 \]
med eksplisitt formel
\[ a_i = 2\cdot3^{i-1} \]
Ved hjelp av summetegnet \sym{$\sum$} kan rekken vår komprimeres betratelig. Ved å skrive $ \sum\limits_{i=1}^5 $ indikerer vi at $ i $ er en løpende variabel som starter på 1, og deretter øker med 1 opp til 5. Da kan vi skrive rekken som $ \sum\limits_{i=1}^5 2\cdot3^{i-1} $, underforstått at vi skal sette et plusstegn hver gang $ i $ øker med 1:
\[ 2+6+18+54+162=\sum\limits_{i=1}^5 2\cdot3^{i-1} \]
Den uendelige rekken 
$ 2+6+18+... $ kan vi skrive som $\sum\limits_{i=1}^\infty 2\cdot3^{i-1}$.

\reg[\regnregsum \label{regnregsum}]{
For to følger $ \lbrace a_i\rbrace $ og $ \lbrace b_i\rbrace $, og en konstant $ c $, har vi at
\begin{align}
	\sum_{i=j}^{n} \left(a_i+b_i\right) &= \sum_{i=j}^{n} a_i + \sum_{i=j}^{n} b_i \label{sumreg1}\br
	\sum_{i=j}^{n} c a_i &= c\sum_{i=j}^{n} a_i \label{sumreg2}
\end{align}
hvor $ {j, n \in \mathbb{N}} $ og $ {j<n} $.
}
\newpage
\fork{\ref{regnregsum} \regnregsum}{
Ved å skrive ut summen og omrokkere på rekkefølgen av addisjonene, innser vi at
\alg{
	\sum_{i=1}^{n} \left(a_i+b_i\right) &= a_1+b_1 + a_2+b_2 + ... + a_n + b_n \\
	&= a_1+a_2+...+a_n + b_1+b_2+...+b_n \\
	&= \sum_{i=1}^{n} a_i + \sum_{i=1}^{n} b_i 
}
Ved å skrive ut summen og faktorisere ut $ c $, innser vi også at
\alg{
	\sum_{i=1}^{n} c a_i &= ca_1 + ca_2 + ...+ ca_n \\
	&= c(a_1+a_2+...+a_n) \\
	&= c\sum_{i=1}^{n} a_i 
}
}
\newpage
\section{Induksjon} \index{induksjon}
I teoretisk matematikk stilles det strenge krav til bevis av formler. En metode som brukes spesielt for formler med heltall, er \outl{induksjon}. Det kan være litt vanskelig i starten å få helt grep på induksjonsprinsippet, så la oss gå rett til et eksempel:\vsk

Vi ønsker å vise at ligningen
\begin{equation}
2+4+6+...+2n=n(n+1) \label{induk}
\end{equation}
er gyldig for alle $n\in\mathbb{N}$. Vi starter med å vise at dette stemmer for $ {n=1} $:
\alg{2 &= 1\cdot(1+1) \\
	2&= 2}
Nå vet vi altså om et heltall, nemlig $ {n=1} $, som ligningen stemmer for. Videre antar vi at ligningen er gyldig helt opp til $n=k $. Vi ønsker å sjekke om den da er gyldig også når $ n=k+1 $. I så fall skal vi ha at
\[ 2+4+6+...+\quad\mathclap{\overbrace{2k}^{\text{ledd nr. }k} \;\;\;\,}+\quad\;\;\;\mathclap{\underbrace{2(k+1)}_{\text{ledd nr. }k+1}}\qquad=(k+1)((k+1)+1) \]
Men frem til ledd nr. $ k $ er det tatt for gitt at \eqref{induk} gjelder, dermed får vi at\footnote{I påfølgende eksempler skal vi for enkelthets skyld la ledd nr. $ k $ være innbakt i symbolet ''$ ... $''. }
\alg{
	\underbrace{2+4+6+...+2k}_{k(k+1)}+2(k+1)&=(k+1)((k+1)+1) \\ 
	k(k+1)+2(k+1)&=(k+1)(k+2)\br
	(k+1)(k+2)&= (k+1)(k+2)
	}
Altså er ligningen gyldig. Og nå kommer den briljante konklusjonen: Vi har vist at \eqref{induk} er sann for ${ n=1 }$. I tillegg har vi vist at hvis ligningen gjelder for ${n= k} $, så gjelder den også for $ {n=k+1} $. På grunn av dette vet vi at \eqref{induk} gjelder for ${n= 1+1=2} $. Men når vi vet at den gjelder for $ {n=2} $, gelder den også for $ {n=2+1=3} $ og så videre, altså for alle heltall!\regv
\reg[Induksjon \label{reginduksjon}]{
	Når vi ved induksjon ønsker å vise at ligningen
	\begin{equation}\label{induksjon}
		A(n)=B(n)
	\end{equation}
	er sann for alle $ n\in\mathbb{N} $, gjør vi følgende:
	\begin{enumerate}
		\item Sjekker at \eqref{induksjon} er sann for $ {n=1} $.
		\item Sjekker at \eqref{induksjon} er sann for $ {n=k+1} $, antatt at den er sann for $ {n=k} $.
	\end{enumerate}
}
\eks[1]{
Vis ved induksjon at
\[ 1+3+5+...+(2n-1) = n^2 \]	
for alle $ n\in\mathbb{N} $.

\sv
Vi sjekker at påstanden stemmer for $ {n=1} $:
\alg{
	1 &= 1^2 \\
	1 &= 1
	}
Vi tar det for gitt at påstanden gjelder for $ {n=k} $, og sjekker at den stemmer også for $ {n=k+1 }$:
\alg{
	\underbrace{1+3+5+...}_{k^2}+(2(k+1)-1) &=(k+1)^2 \\
	k^2 + 2k+1 &= (k+1)^2 \br
	(k+1)^2 &= (k+1)^2
	}
Dermed er påstanden vist for alle $ n\in\mathbb{N} $.\vsk
}
\newpage
\eks[2]{
Vis ved induksjon at
\[ 1^3 + 2^3 + 3^3 + ... + n^3= \frac{n^2(n+1)^2}{4} \]
for alle $ n\in\mathbb{N} $. \\

\sv
Vi starter med å sjekke for $ n=1 $:
\alg{1 &= \frac{1^2\cdot(1+1)^2}{4} \\
	1^3&= \frac{2^2}{4} \\
	1 &= 1}
Ligningen er altså sann for $ {n=1 }$. Vi antar videre at den også stemmer for $ {n=k} $, og sjekker for $ {n=k+1} $:
\alg{\underbrace{1^3+2^3+3^3+...}_{\frac{k^2(k+1)^2}{4}}+(k+1)^3 &= \frac{(k+1)^2(k+1+1)^2}{4} \\
	\frac{k^2(k+1)^2}{4} + (k+1)^3 &= \frac{(k+1)^2(k+2)^2}{4} \\
	\frac{k^2(k+1)^2+4(k+1)^3}{4} &= \\
	\frac{(k+1)^2(k^2+4(k+1))}{4} &= \\	
	\frac{(k+1)^2(k^2+4k+4))}{4} &= \\	
	\frac{(k+1)^2(k+2)^2}{4} &= \frac{(k+1)^2(k+2)^2}{4}
}
Påstanden er dermed vist for alle $ n\in\mathbb{N} $.		
	}\newpage
\eks[3]{\label{prodind}
Vis ved induksjon at
\[ 3\cdot9\cdot27\cdot... \cdot 3^n= 3^{\frac{1}{2}n(n+1)}\]
for alle $ n\in\mathbb{N} $.

\sv
Vi sjekker at påstanden er sann for $ n=1 $:
\alg{
3 &= 3^{\frac{1}{2}\cdot1(1+1)} \\
3 &= 3^1
}
Videre antar vi at påstanden stemmer også for $ {n=k}$, og sjekker for $ n=k+1 $:
\alg{
 \underbrace{3\cdot9\cdot27\cdot...}_{3^{\frac{1}{2}k(k+1)}} \cdot 3^{k+1} &= 3^{\frac{1}{2}(k+1)(k+1+1)} \br
3^{\frac{1}{2}k(k+1)}\cdot  3^{k+1} &= 3^{\frac{1}{2}(k+1)(k+2)} \\
3^{\frac{1}{2}k(k+1)+k+1} &= \\
3^{\frac{1}{2}k(k+1)+\frac{2}{2}(k+1)} &= \\
3^{\frac{1}{2}(k+1)(k+2)} &= 3^{\frac{1}{2}(k+1)(k+2)}
}
Påstanden er dermed vist for alle $ n\in \mathbb{N} $.
}
\info{Tips}{
	Det kan oppstå tilfeller hvor det å faktorisere uttrykk kan være vanskelig. Da kan man i stedet skrive ut alle uttrykk for å undersøke om ligningen er gyldig.
}

\end{document}
